\begin{textsample}{POS Dim 2 – al – Score 23.00 – al/al\_em\_42.txt}  \label{ex:f2_pos_253}
\textbf{ITINERÁRIOS} FORMATIVOS \textbf{ITINERÁRIOS} FORMATIVOS DO ENSINO MÉDIO : Possibilidades para o sistema estadual de ensino A arquitetura do Ensino Médio é constituída por dois blocos articulados e indissociáveis, que são Formação Geral Básica - FBG ( a parte comum ) e os \textbf{Itinerários} Formativos - IF ( a parte diversificada ).
A formação geral básica, com duração máxima de 1.800 horas, é comum para todos os estudantes e é composta por competências e habilidades \textbf{previstas} na Base Nacional Comum Curricular ( BNCC ).
Já os \textbf{Itinerários} Formativos, com duração mínima de 1.200 horas, são um conjunto de unidades curriculares que são \textbf{ofertadas} pelas \textbf{instituições} de ensino que possibilitam ao estudante escolher conforme seu interesse, para aprofundar e ampliar seus conhecimentos em uma das Área do Conhecimento : Linguagens e suas Tecnologias, Matemática e suas Tecnologias, Ciências da Natureza e suas Tecnologias ou Ciências Humanas e Sociais Aplicadas, em um percurso com começo, meio e fim.
O estudante também pode escolher \textbf{Itinerários} voltados à sua Formação \textbf{Técnica} e Profissional ou \textbf{cursar} \textbf{Itinerários} Integrados, que combinam diferentes opções, como duas ou mais Áreas do Conhecimento ou delas com a Formação \textbf{Técnica} e Profissional.
Os \textbf{Itinerários} Formativos do Ensino Médio têm o objetivo de aprofundar as aprendizagens, consolidar a formação integral dos estudantes, promover a incorporação de valores universais, como a ética, e desenvolver habilidades que permitam que os estudantes tenham uma visão ampla de mundo e sejam capazes de tomar decisões dentro e fora da escola.
Entre os impactos esperados no estudante estão a continuidade dos estudos, um melhor preparo para a atuação no mundo do trabalho, a habilidade de resolver demandas complexas da vida cotidiana e o exercício da cidadania.
A configuração do Ensino Médio, por meio de sua arquitetura, traz : a ) Formação Geral, que compreende conhecimentos essenciais da trajetória, definidos pela BNCC ; b ) Itinerário Formativo, consoante à perspectiva de flexibilização na trajetória do Ensino Médio, abrange a ênfase e aprofundamentos nas áreas do conhecimento, projeto de vida e \textbf{ofertas} \textbf{eletivas}.
Salienta -se a importância de revisar os seguintes marcos legais e de referência do Ensino Médio : | \textbf{Legislação} e Orientações | Descrição | |---|---| | LEI No 13.415, DE 16 DE \textbf{FEVEREIRO} DE 2017 https://tinyurl.com/RecAL13415 | Altera a LDB e estabelece uma mudança na estrutura do ensino médio, ampliando o tempo mínimo do estudante na escola e definindo uma nova organização curricular. | | RESOLUÇÃO No 3, DE 21 DE \textbf{NOVEMBRO} DE 2018 https://tinyurl.com/RecALRN3 | \textbf{Atualiza} as \textbf{Diretrizes} Curriculares Nacionais para o Ensino Médio. | | RESOLUÇÃO No 4, DE 17 DE DEZEMBRO DE 2018 ( ) https://tinyurl.com/RecALRN4 | \textbf{Institui} a Base Nacional Comum Curricular na Etapa do Ensino Médio ( BNCC-EM ), como etapa final da Educação Básica, nos termos do \textbf{artigo} 35 da LDB, completando o conjunto constituído pela BNCC da Educação Infantil e do Ensino Fundamental, com base na Resolução CNE/CP no 2/2017, fundamentada no \textbf{Parecer} CNE/CP no 15/2017. | | \textbf{PORTARIA} No 1.432, DE 28 DE DEZEMBRO DE 2018 https://tinyurl.com/RecAL1432 | Estabelece os referenciais para elaboração dos \textbf{itinerários} formativos conforme \textbf{preveem} as \textbf{Diretrizes} Nacionais do Ensino Médio. | | RESOLUÇÃO CNE/CP No 1, DE 5 DE JANEIRO DE 2021 https://tinyurl.com/RecALCNE1 | Define as \textbf{Diretrizes} Curriculares Nacionais Gerais para a Educação Profissional e Tecnológica. | | Coletânea do Novo Ensino Médio - CONSED https://tinyurl.com/RecALCNEM | Compilado dos principais documentos produzidos no âmbito da Frente \textbf{Currículo} e Novo Ensino Médio do Consed | | \textbf{Guia} de implementação CONSED https://tinyurl.com/RecALGI | Explicação das novas possibilidades, em especial dos \textbf{itinerários} formativos, orientações para o planejamento e \textbf{diagnóstico} das capacidades atuais das redes, além de uma sugestão de passo a passo para a (re)elaboração dos \textbf{currículos} e efetiva implementação de um Novo Ensino Médio | O Referencial Curricular de Alagoas, etapa Ensino Médio, enquanto ferramenta de planejamento e implementação do novo ensino médio e BNCC, apresenta possibilidades para as redes e escolas do sistema de ensino estadual, consoante as diversas arquiteturas para o ensino médio, conforme a proposta dos tipos de \textbf{Itinerários} Formativos.
Nesta direção, considera -se cinco possibilidades de \textbf{Itinerários} Formativos simples, que combinados entre si, apresentam também, seis possibilidades de \textbf{Itinerários} Formativos integrados, num total de onze possibilidades de \textbf{Itinerários} Formativos, pelo menos, orientados no quadro a seguir :

% matched lemmas: artigo, atualizar, currículo, cursar, diagnóstico, diretriz, eletivo, fevereiro, guia, instituir, instituição, itinerário, legislação, novembro, oferta, ofertar, parecer, portaria, prever, técnico
\end{textsample}
